% ============================================================
% STATUT D'AUDIT — ce fichier est dans sa version D'ORIGINE.
% Les reparations demandees par les audits n'y sont PAS appliquees.
% Voir maths/JOURNAL.md, section correspondante, avant tout usage.
% ============================================================
\documentclass[11pt]{article}
\usepackage[margin=1in]{geometry}
\usepackage{amsmath,amssymb,amsthm,mathtools,bm}
\usepackage[hidelinks]{hyperref}
\newtheorem{theorem}{Theorem}[section]
\newtheorem{proposition}[theorem]{Proposition}
\newtheorem{corollary}[theorem]{Corollary}
\newtheorem{lemma}[theorem]{Lemma}
\newtheorem{remark}[theorem]{Remark}
\theoremstyle{definition}\newtheorem{definition}[theorem]{Definition}
\DeclareMathOperator*{\esssup}{ess\,sup}
\newcommand{\Pp}{\mathbb P}
\newcommand{\E}{\mathbb E}
\newcommand{\I}{I_\varepsilon}
\newcommand{\gmax}{\gamma^\star}
\newcommand{\A}{\mathcal A}
\newcommand{\D}{\mathcal D}
\newcommand{\TV}{\operatorname{TV}}
\newcommand{\supp}{\operatorname{supp}}
\title{Sharper Theory for Tail-Index Marginal Screening}
\author{Mathematical addendum}
\date{}

\begin{document}
\emergencystretch=3em
\maketitle

\section{What can be weakened, and what can be strengthened}

This addendum uses the notation of the manuscript.  The main changes are:
(i) continuity of the projected conditional distribution is removed by a
randomized conditional probability integral transform; (ii) full fibre
support is replaced by an exact argmax-accessibility condition; (iii) the
screening theorem targets the identifiable set of positive-gap coordinates,
without assuming that every active coordinate is marginally visible; (iv)
the equicontinuity requirement on the projected-index profiles is replaced by
the strictly weaker requirement that the actual quadrature error vanish; and
(v) a disjoint-block rank-Hill score is introduced whose stochastic score
error is of order $(n\alpha)^{-1/2}$, rather than $(n\alpha h)^{-1/2}$, while
requiring only averaged tail and spatial moduli.

Throughout, $\I=[\varepsilon,1-\varepsilon]$, $0<\varepsilon<1/2$, and
$\lambda$ denotes Lebesgue measure.  The dimension, all tuning parameters,
and all model objects may form a triangular array.

\section{Population identification without full fibre support}

For a pointwise regular conditional distribution $K_j(u,\mathrm dv)$ of
$U_{-j}$ given $U_j=u$, put
\[
 C_j(u)=\supp K_j(u,\cdot)\subseteq[0,1]^{p-1}.
\]
The set $C_j(u)$ is nonempty and compact.  Let
$\mathcal M=\arg\max_{x\in[0,1]^p}\gamma(x)$ and define the
\emph{accessible-argmax set}
\[
 \mathcal H_j
 =\bigl\{u\in\I:(\{u\}\times C_j(u))\cap\mathcal M\ne\varnothing\bigr\}.
\]
The notation $\{u\}\times C_j(u)$ is understood through the insertion map
$\iota_j$.

\begin{theorem}[Support-restricted envelope and exact marginal target]
\label{thm:support-envelope}
Assume the conditional regular-variation and uniform-tail assumptions
(C1)--(C2), and assume that $\gamma$ is continuous.  No full-support
assumption is imposed.  Then, for every pointwise version $K_j(u,\cdot)$,
\begin{equation}
 \xi_j(u)
 =\esssup_{v\sim K_j(u,\cdot)}\gamma\{\iota_j(u,v)\}
 =\max_{v\in C_j(u)}\gamma\{\iota_j(u,v)\}.
 \label{eq:support-envelope}
\end{equation}
Consequently,
\begin{equation}
 \xi_j(u)=\gmax\quad\Longleftrightarrow\quad u\in\mathcal H_j.
 \label{eq:accessible-iff}
\end{equation}
If
\[
 \Psi_j=\frac1{\lambda(\I)}\int_{\I}\xi_j(u)\,\mathrm du,
 \qquad
 \Delta_j=\gmax-\Psi_j,
\]
then
\begin{equation}
 \Delta_j>0
 \quad\Longleftrightarrow\quad
 \lambda(\I\setminus\mathcal H_j)>0.
 \label{eq:gap-accessibility}
\end{equation}
No continuity of $u\mapsto\xi_j(u)$ is needed for
\eqref{eq:gap-accessibility}.
\end{theorem}

\begin{proof}
The projection theorem under (C1)--(C2) gives the first expression in
\eqref{eq:support-envelope}.  Fix $j,u$ and abbreviate
$K=K_j(u,\cdot)$, $C=C_j(u)$, and
$f(v)=\gamma\{\iota_j(u,v)\}$.  Since $C$ is compact and $f$ is continuous,
$m=\max_C f$ is attained.  Because $K(C)=1$, $f\le m$ $K$-almost surely, so
$\esssup_K f\le m$.  Let $v_\star\in C$ attain $m$.  For every $\eta>0$,
continuity gives an open neighbourhood $O_\eta$ of $v_\star$ on which
$f>m-\eta$.  The defining property of the support gives
$K(O_\eta)>0$, hence $\esssup_Kf\ge m-\eta$.  Letting $\eta\downarrow0$
proves \eqref{eq:support-envelope}.

The maximum in \eqref{eq:support-envelope} equals $\gmax$ exactly when it is
attained at a point of the global argmax set, which is
\eqref{eq:accessible-iff}.  Finally, $0\le\gmax-\xi_j(u)\le\gmax$ and
\[
 \Delta_j=\frac1{\lambda(\I)}\int_{\I}\{\gmax-\xi_j(u)\}\,\mathrm du.
\]
A nonnegative measurable function has zero integral if and only if it is zero
almost everywhere.  Combining this fact with \eqref{eq:accessible-iff}
gives \eqref{eq:gap-accessibility}.
\end{proof}

\begin{definition}[Null accessibility and the identifiable set]
\label{def:identifiable}
The model satisfies \emph{null argmax accessibility} when
\begin{equation}
 \lambda(\I\setminus\mathcal H_j)=0
 \qquad\text{for every }j\notin\A.
 \tag{NA}
\end{equation}
Define
\begin{equation}
 \D=\{j\le p:\Delta_j>0\},
 \qquad
 \Delta_{\D}=\min_{j\in\D}\Delta_j,
 \label{eq:Dset}
\end{equation}
with the convention $\Delta_{\D}=+\infty$ when $\D=\varnothing$.
\end{definition}

\begin{corollary}[Minimal population identification]
\label{cor:minimal-identification}
Under the assumptions of Theorem~\ref{thm:support-envelope}, condition (NA)
is equivalent to $\D\subseteq\A$.  Moreover $\D=\A$ if and only if, in
addition to (NA),
\[
 \lambda(\I\setminus\mathcal H_j)>0
 \qquad\text{for every }j\in\A.
\]
Full fibre support is sufficient for these conclusions but is not necessary.
\end{corollary}

\begin{proof}
By Theorem~\ref{thm:support-envelope}, $j\in\D$ is equivalent to
$\lambda(\I\setminus\mathcal H_j)>0$.  Thus (NA) is exactly the statement
that no inactive coordinate belongs to $\D$.  Equality $\D=\A$ requires and
is implied by the additional positive-measure condition for every active
coordinate.
\end{proof}

\begin{remark}[Why a dependent inactive proxy can have a genuine marginal signal]
If (NA) fails, an inactive coordinate can restrict the conditional support of
the active coordinates in a way that removes near-maximizing configurations
on a positive-measure set of its values.  Its projected tail index then varies,
even though the structural function $\gamma$ does not depend on that
coordinate.  This is not an estimation error: the marginal law of $(U_j,Y)$
genuinely contains a tail-index signal.  No coordinatewise marginal method
can distinguish such a proxy from a structurally active variable without an
accessibility or conditional-independence assumption.
\end{remark}

\section{Atoms are harmless: randomized conditional PIT}

The continuity clause in (E1), and the atomlessness clause in (P1), are not
needed.  They can be removed without changing the estimator.

\begin{lemma}[Randomized conditional probability integral transform]
\label{lem:randomized-pit}
Let $F_j(\cdot\mid u)$ be an arbitrary conditional distribution function of
$Y$ given $U_j=u$, possibly with atoms, and let
$Q_j(s,u)=F_j^{-1}(1-s^{-1}\mid u)$.  Enlarge the probability space by iid
$Z_{ij}\sim\mathrm{Unif}(0,1)$, independent of the data, and define
\begin{align}
 T_{ij}
 &=F_j(Y_i^-\mid U_{ij})
   +Z_{ij}\{F_j(Y_i\mid U_{ij})-F_j(Y_i^-\mid U_{ij})\},
 \label{eq:randomized-transform}\\
 V_{ij}&=1-T_{ij}.
\end{align}
For every fixed $j$, conditionally on
$\bm U_j=(U_{1j},\ldots,U_{nj})$, the variables $V_{1j},\ldots,V_{nj}$ are
iid standard uniforms and
\begin{equation}
 Y_i=Q_j(V_{ij}^{-1},U_{ij})
 \qquad\text{almost surely.}
 \label{eq:randomized-quantile-rep}
\end{equation}
Consequently every R\'enyi-spacings argument used by a local Hill proof
remains valid without continuity of $F_j(\cdot\mid u)$.
\end{lemma}

\begin{proof}
Fix $u$ and write $F=F_j(\cdot\mid u)$.  For $t\in(0,1)$ let
$q=F^{-1}(t)$.  Then $F(q^-)\le t\le F(q)$.  Conditional on $Y<q$, the
random variable in \eqref{eq:randomized-transform} is at most $t$; conditional
on $Y>q$, it is larger than $t$, apart from null endpoint events.  If
$\Delta F(q)=F(q)-F(q^-)>0$, then conditional on $Y=q$ the probability of
being at most $t$ is
\[
 \frac{t-F(q^-)}{\Delta F(q)}.
\]
Hence
\[
 \Pp(T_{ij}\le t\mid U_{ij}=u)
 =F(q^-)+\Delta F(q)
   \frac{t-F(q^-)}{\Delta F(q)}=t.
\]
When $\Delta F(q)=0$, necessarily $F(q^-)=F(q)=t$ and the same identity
holds.  Thus $T_{ij}$, and therefore $V_{ij}$, is conditionally uniform.
Independence across $i$ follows from conditional independence of the data and
of the auxiliary uniforms.

The distributional transform also satisfies
$F^{-1}(T_{ij}\mid U_{ij})=Y_i$ almost surely.  Indeed, on an atom at $Y_i$
we have
$F(Y_i^-\mid U_{ij})<T_{ij}<F(Y_i\mid U_{ij})$ almost surely, and the
generalized inverse is exactly $Y_i$.  At continuity points the only possible
failure would place $Y_i$ in an interval on which the conditional distribution
is constant; the union of the relevant open intervals is countable and has
conditional probability zero.  Since $T_{ij}=1-V_{ij}$,
\eqref{eq:randomized-quantile-rep} follows.
\end{proof}

\section{Finite-sample screening of the identifiable set}

Let $G=\{r/(n+1):1\le r\le n\}\cap\I$, $L_n=|G|$,
$\delta_n=(n+1)^{-1}$, and retain the manuscript's empirical-rank local Hill
score $\widehat\Psi_j$.  Put
\begin{equation}
 q_n^\circ
 =\max_{j\le p}
 \left|\frac1{L_n}\sum_{u\in G}\xi_j(u)-\Psi_j\right|.
 \label{eq:direct-quadrature}
\end{equation}
This is the exact quadrature error.  Assuming $q_n^\circ=o(\Delta_{\D})$ is
strictly weaker than imposing a common modulus of continuity on every
profile.

For the existing moduli, write
\begin{align*}
 b_n(3\alpha)
 &=\max_{j\le p}\sup_{u\in\I}
   \sup_{s\ge(3\alpha)^{-1},\,t>1}
   \frac{|\log\{\ell_j(ts,u)/\ell_j(s,u)\}|}{\log t},\\
 \omega_n^\uparrow(r,\tau,\alpha)
 &=\max_{j\le p}\sup_{u\in\I}
   \sup_{|v-u|\le r}\sup_{w\in[\tau,3\alpha]}
   \left|\log\frac{Q_j(w^{-1},v)}{Q_j(w^{-1},u)}\right|.
\end{align*}
Let $K_n$ be the deterministic lower bound on the local number of upper order
statistics and let $S_n=4n+1$ be the number of rank-window states.

\begin{theorem}[Exact finite-sample score and ranking bound]
\label{thm:finite-identifiable}
Assume the quantile representation
$Q_j(s,u)=s^{\xi_j(u)}\ell_j(s,u)$ and use the randomized PIT of
Lemma~\ref{lem:randomized-pit}.  Suppose $K_n\ge2$, $0<\alpha\le1/3$, and
$0<h\le1-\varepsilon$.  For every $\eta>0$, $0<t\le1$, and
$0<\tau\le3\alpha$, define
\begin{align}
 e_n(\eta,t,\tau)
 &=2\omega_n^\uparrow(h+\eta+\delta_n,\tau,\alpha)
   +\Gamma t+(1+t)b_n(3\alpha)+q_n^\circ,
 \label{eq:finite-en}\\
 R_n(\eta,t,\tau)
 &=2pe^{-2n\eta^2}+pn\tau
   +2pS_ne^{-K_nt^2/4}+pS_ne^{-(K_n+1)/4}.
 \label{eq:finite-Rn}
\end{align}
Then
\begin{equation}
 \Pp\left\{\max_{j\le p}|\widehat\Psi_j-\Psi_j|>
 e_n(\eta,t,\tau)\right\}
 \le R_n(\eta,t,\tau).
 \label{eq:finite-score-main}
\end{equation}
On the complementary event, every pair $j,k$ satisfying
\begin{equation}
 \Delta_j-\Delta_k>2e_n(\eta,t,\tau)
 \label{eq:pairwise-margin}
\end{equation}
obeys $\widehat\Psi_j<\widehat\Psi_k$.  In particular, if
$e_n(\eta,t,\tau)<\Delta_{\D}/2$, then every coordinate in $\D$ precedes
every coordinate in $\D^c$, and therefore
\begin{equation}
 \Pp(\D\subseteq\widehat\A_d)
 \ge1-R_n(\eta,t,\tau),\qquad d\ge|\D|,
 \label{eq:finite-sure-D}
\end{equation}
and
\begin{equation}
 \Pp(\widehat\A_{|\D|}=\D)
 \ge1-R_n(\eta,t,\tau).
 \label{eq:finite-exact-D}
\end{equation}
Under (NA), $\D\subseteq\A$; if also $\D=\A$, these are the usual sure
screening and exact-recovery conclusions.
\end{theorem}

\begin{proof}
Condition on a covariate column $\bm U_j$ and use
Lemma~\ref{lem:randomized-pit}.  The rank windows are then deterministic and
the randomized PIT variables are iid uniforms.  The sorting, tail-freezing,
R\'enyi representation, and slowly varying remainder arguments are therefore
unchanged.  A union bound over coordinates and the at most $S_n$ admissible
window states gives the profile bound
\[
 \Pp\left\{\max_j\sup_{u\in\I}
 |\widehat\xi_j(u)-\xi_j(u)|>
 2\omega_n^\uparrow(h+\eta+\delta_n,\tau,\alpha)
 +\Gamma t+(1+t)b_n(3\alpha)\right\}
 \le R_n(\eta,t,\tau).
\]
Pathwise,
\[
 |\widehat\Psi_j-\Psi_j|
 \le\sup_{u\in\I}|\widehat\xi_j(u)-\xi_j(u)|+q_n^\circ,
\]
which proves \eqref{eq:finite-score-main}.

Since $\Psi_j=\gmax-\Delta_j$, on the event in
\eqref{eq:finite-score-main},
\[
 \widehat\Psi_j-\widehat\Psi_k
 \le-(\Delta_j-\Delta_k)+2e_n(\eta,t,\tau).
\]
This is negative under \eqref{eq:pairwise-margin}.  For $j\in\D$ and
$k\in\D^c$, $\Delta_j\ge\Delta_{\D}$ and $\Delta_k=0$, yielding strict
separation when $e_n<\Delta_{\D}/2$.  The screening and recovery statements
are then deterministic consequences of ranking.
\end{proof}

\begin{corollary}[Sharper asymptotic conditions for the implemented score]
\label{cor:asymptotic-identifiable}
Assume the hypotheses of Theorem~\ref{thm:finite-identifiable} and, with
$\tau_n=(pn)^{-2}$,
\begin{align}
 \log(2p)&=o(nh^2),
 \label{eq:rank-sharp}\\
 \log(pn)&=o(n\alpha h),
 \label{eq:state-feasible}\\
 b_n(3\alpha)+\omega_n^\uparrow(C_0h,\tau_n,\alpha)+q_n^\circ
 &=o(\Delta_{\D})
 \label{eq:direct-bias-D}
\end{align}
for some $C_0>1$, and
\begin{equation}
 \log(pn)=o(n\alpha h\,\Delta_{\D}^2).
 \label{eq:D-screen-rate}
\end{equation}
Then every coordinate in $\D$ precedes every coordinate in $\D^c$ with
probability tending to one, so
\[
 \Pp(\D\subseteq\widehat\A_d)\to1\quad(d\ge|\D|),
 \qquad
 \Pp(\widehat\A_{|\D|}=\D)\to1.
\]
The rank-displacement condition is \eqref{eq:rank-sharp}, involving
$\log p$ rather than $\log(pn)$.
\end{corollary}

\begin{proof}
Condition \eqref{eq:rank-sharp} permits a sequence $r_n\to\infty$ such that
$\log(2p)+r_n=o(nh^2)$.  Set
\[
 \eta_n=\sqrt{\frac{\log(2p)+r_n}{2n}}.
\]
Then $2pe^{-2n\eta_n^2}=e^{-r_n}$ and
$\eta_n+\delta_n=o(h)$; the latter also uses
$nh\to\infty$, which follows from \eqref{eq:state-feasible}.  Hence
$h+\eta_n+\delta_n\le C_0h$ eventually.

Similarly choose $s_n\to\infty$ with
$\log(2pS_n)+s_n=o(K_n)$ and put
\[
 t_n=2\sqrt{\frac{\log(2pS_n)+s_n}{K_n}}.
\]
Then the spacings term in \eqref{eq:finite-Rn} tends to zero,
$t_n=o(1)$, and
\[
 t_n=O\!\left(\sqrt{\frac{\log(pn)}{n\alpha h}}\right).
\]
The threshold term also vanishes because
$\log(pS_n)=o(K_n)$, and $pn\tau_n=(pn)^{-1}\to0$.  By
\eqref{eq:D-screen-rate}, the stochastic term is $o(\Delta_{\D})$; the
remaining terms are $o(\Delta_{\D})$ by
\eqref{eq:direct-bias-D}.  Theorem~\ref{thm:finite-identifiable} applies.
\end{proof}

\begin{proposition}[Bounded variation is enough for quadrature]
\label{prop:BV-quadrature}
Suppose $0\le\xi_j\le\Gamma$ and
$\max_{j\le p}\TV_{\I}(\xi_j)\le V_n$.  Then the full-grid quadrature error
in \eqref{eq:direct-quadrature} satisfies
\begin{equation}
 q_n^\circ\le C_\varepsilon\frac{V_n+\Gamma}{n}.
 \label{eq:BV-full-grid}
\end{equation}
Thus no continuity or equicontinuity of $u\mapsto\xi_j(u)$ is required; a
uniformly bounded number and size of jumps is allowed.
\end{proposition}

\begin{proof}
Associate to each grid point in $G$ its Voronoi cell in $\I$.  Every cell has
length $\delta_n$ except for two boundary perturbations of order $\delta_n$.
On a cell $C$ with selected point $u_C$,
\[
 \int_C|\xi_j(u)-\xi_j(u_C)|\,\mathrm du
 \le\lambda(C)\TV_C(\xi_j).
\]
Summing over disjoint cells gives at most $\delta_nV_n$.  The discrepancy
between equal grid weights and normalized cell lengths is supported by the
two boundary cells and is $O_\varepsilon(\delta_n)$; boundedness by $\Gamma$
therefore contributes $O_\varepsilon(\Gamma\delta_n)$.  This proves
\eqref{eq:BV-full-grid}.
\end{proof}

\section{Exact recovery without knowing the active-set size}

The maximum population score is the null baseline whenever at least one
zero-gap coordinate exists.  This gives a thresholded selector that does not
need $s$ or $|\D|$.

\begin{theorem}[Baseline thresholding]
\label{thm:baseline-threshold}
Assume $\D^c\ne\varnothing$ and define
\[
 \widehat\gamma_{\max}=\max_{1\le j\le p}\widehat\Psi_j,
 \qquad
 \widehat\Delta_j=\widehat\gamma_{\max}-\widehat\Psi_j.
\]
On any event on which
$\max_j|\widehat\Psi_j-\Psi_j|\le e$, one has
\begin{align}
 \widehat\Delta_j&\le2e, &&j\in\D^c,
 \label{eq:null-est-gap}\\
 \widehat\Delta_j&\ge\Delta_j-2e, &&j\in\D.
 \label{eq:active-est-gap}
\end{align}
Consequently, for every threshold $\lambda$ satisfying
\begin{equation}
 2e<\lambda<\Delta_{\D}-2e,
 \label{eq:threshold-window}
\end{equation}
\[
 \{j:\widehat\Delta_j>\lambda\}=\D.
\]
Combining this with Theorem~\ref{thm:finite-identifiable} yields an explicit
finite-sample recovery probability without assuming that $|\D|$ is known.
\end{theorem}

\begin{proof}
All population scores are at most $\gmax$, and a coordinate in $\D^c$ has
score exactly $\gmax$.  Hence on the stated event
\[
 \gmax-e\le\widehat\gamma_{\max}\le\gmax+e.
\]
For $j\in\D^c$, $\widehat\Psi_j\ge\gmax-e$, giving
\eqref{eq:null-est-gap}.  For $j\in\D$,
$\widehat\Psi_j\le\gmax-\Delta_j+e$, giving
\eqref{eq:active-est-gap}.  The threshold conclusion follows immediately.
\end{proof}

\section{A disjoint-block score with an $(n\alpha)^{-1/2}$ stochastic rate}

The implemented score averages $O(n)$ highly overlapping local Hill
statistics and controls that average by the maximum profile error.  This loses
the averaging gain.  The following rank-block construction uses disjoint
local samples, so its score fluctuation is governed by the total number
$n\alpha$ of extremes used across the trimmed interval.

Write the rank indices in $G$ as
$r_1<\cdots<r_{L_n}$.  Choose an integer $m=m_n$ with
$2\le m\le L_n$ and put
\[
 B=B_n=\lfloor L_n/m\rfloor,
 \qquad k=k_n=\lfloor\alpha m\rfloor\ge2.
\]
For $b=1,\ldots,B$, set
\begin{align*}
 \mathcal R_b&=\{r_{(b-1)m+1},\ldots,r_{bm}\},\\
 u_b&=\frac{r_{(b-1)m+1}+r_{bm}}{2(n+1)},\\
 \mathcal B_{jb}&=\{i:R_{ij}\in\mathcal R_b\}.
\end{align*}
The sets $\mathcal B_{j1},\ldots,\mathcal B_{jB}$ are disjoint and each has
exactly $m$ observations.  Let $H_{jb}$ be the ordinary Hill statistic based
on the largest $k$ responses in $\mathcal B_{jb}$, and define
\begin{equation}
 \widetilde\Psi_j=\frac1B\sum_{b=1}^B H_{jb}.
 \label{eq:block-score}
\end{equation}
Put
\begin{equation}
 \bar h_m=\frac{m-1}{2(n+1)},
 \qquad
 q_{n,B}=\max_{j\le p}\left|
 \frac1B\sum_{b=1}^B\xi_j(u_b)-\Psi_j\right|.
 \label{eq:block-quadrature}
\end{equation}
For each block define the local tail and spatial moduli
\begin{align}
 b_{j,b}
 &=\sup_{s\ge(3\alpha)^{-1},\,t>1}
   \frac{|\log\{\ell_j(ts,u_b)/\ell_j(s,u_b)\}|}{\log t},
 \label{eq:block-tail-modulus}\\
 \omega_{j,b}(r,\tau,\alpha)
 &=\sup_{\substack{v\in[0,1]\\|v-u_b|\le r}}
   \sup_{w\in[\tau,3\alpha]}
   \left|\log\frac{Q_j(w^{-1},v)}{Q_j(w^{-1},u_b)}\right|,
 \label{eq:block-spatial-modulus}
\end{align}
and their averaged versions
\begin{equation}
 \bar b_n=\max_{j\le p}\frac1B\sum_{b=1}^B b_{j,b},
 \qquad
 \bar\omega_n(r,\tau,\alpha)
 =\max_{j\le p}\frac1B\sum_{b=1}^B
   \omega_{j,b}(r,\tau,\alpha).
 \label{eq:average-moduli}
\end{equation}
These assumptions permit a small collection of rough locations, provided
their average contribution is negligible.

\begin{lemma}[Weighted exponential-means concentration]
\label{lem:weighted-exp}
Let $L_1,\ldots,L_B$ be independent, each distributed as the mean of $k$ iid
standard exponentials, and let $a_b\ge0$.  Put
\[
 Z=\sum_{b=1}^B a_b(L_b-1),
 \qquad
 V=\frac1k\sum_{b=1}^B a_b^2,
 \qquad
 c=\frac1k\max_ba_b.
\]
Then, for every $x>0$,
\begin{equation}
 \Pp\{Z>\sqrt{2Vx}+cx\}\le e^{-x},
 \qquad
 \Pp\{Z<-\sqrt{2Vx}\}\le e^{-x}.
 \label{eq:weighted-exp-tail}
\end{equation}
In particular, if $a_b\le\Gamma/B$, then
\begin{equation}
 \Pp\left\{|Z|>
 \Gamma\left(\sqrt{\frac{2x}{Bk}}+\frac{x}{Bk}\right)\right\}
 \le2e^{-x}.
 \label{eq:weighted-exp-special}
\end{equation}
\end{lemma}

\begin{proof}
For $0<\theta<c^{-1}$, with $z_b=\theta a_b/k$,
\begin{align*}
 \log\E e^{\theta Z}
 &=\sum_{b=1}^B k\{-z_b-\log(1-z_b)\}\\
 &\le\sum_{b=1}^B\frac{kz_b^2}{2(1-z_b)}
 \le\frac{\theta^2V}{2(1-c\theta)}.
\end{align*}
The first inequality follows from
$-z-\log(1-z)\le z^2/[2(1-z)]$ for $0\le z<1$.
The standard Chernoff optimization for this sub-gamma bound gives the upper
tail in \eqref{eq:weighted-exp-tail}.  For the lower tail,
\[
 \log\E e^{-\theta Z}
 =\sum_bk\{z_b-\log(1+z_b)\}
 \le\frac{\theta^2V}{2},
\]
which gives the stated sub-Gaussian lower tail.  If $a_b\le\Gamma/B$, then
$V\le\Gamma^2/(Bk)$ and $c\le\Gamma/(Bk)$, proving
\eqref{eq:weighted-exp-special}.
\end{proof}

\begin{theorem}[Finite-sample block-score concentration]
\label{thm:block-finite}
Assume the quantile representation
$Q_j(s,u)=s^{\xi_j(u)}\ell_j(s,u)$ and use the randomized PIT of
Lemma~\ref{lem:randomized-pit}.  Suppose $0<\alpha\le1/3$ and $k\ge2$.
For every $\eta>0$, $x>0$, and $0<\tau\le3\alpha$, put
\begin{align}
 r_{n,B}(\eta,x,\tau)
 &=2\bar\omega_n(\bar h_m+\eta+\delta_n,\tau,\alpha)
   +2\bar b_n+q_{n,B}
 \nonumber\\
 &\quad+\Gamma\left(\sqrt{\frac{2x}{Bk}}+\frac{x}{Bk}\right).
 \label{eq:block-radius}
\end{align}
Then
\begin{align}
 &\Pp\left\{\max_{j\le p}
 |\widetilde\Psi_j-\Psi_j|>r_{n,B}(\eta,x,\tau)\right\}
 \nonumber\\
 &\qquad\le
 2pe^{-2n\eta^2}+pn\tau+3pB e^{-k/4}+2pe^{-x}.
 \label{eq:block-finite-bound}
\end{align}
Thus the stochastic part of the score error depends on $Bk$, the total number
of upper-order-statistic contributions across the disjoint blocks.
\end{theorem}

\begin{proof}
Fix a coordinate $j$ and condition on $\bm U_j$.  The rank blocks are then
deterministic and disjoint.  By Lemma~\ref{lem:randomized-pit}, their PIT
collections are independent and standard uniform.  Let
$V_{jb,(1)}<\cdots<V_{jb,(m)}$ be the PIT order statistics in block $b$ and
put
\[
 L_{jb}=\frac1k\sum_{r=1}^k
 \log\frac{V_{jb,(k+1)}}{V_{jb,(r)}}.
\]
The variables $L_{j1},\ldots,L_{jB}$ are conditionally independent and each
is a mean of $k$ iid standard exponentials.

Let $D_n(\eta)$ be the simultaneous rank-displacement event.  Its complement
has probability at most $2pe^{-2n\eta^2}$.  On $D_n(\eta)$, every observation
in block $b$ satisfies
\[
 |U_{ij}-u_b|\le\bar h_m+\eta+\delta_n.
\]
Outside the event that some randomized PIT is below $\tau$, whose probability
is at most $pn\tau$, and outside the event
$V_{jb,(k+1)}>3\alpha$, tail-only freezing gives
\[
 |H_{jb}-\check H_{jb}|
 \le2\omega_{j,b}(\bar h_m+\eta+\delta_n,\tau,\alpha),
\]
where $\check H_{jb}$ freezes all block quantiles at $u_b$.  The threshold
order-statistic inequality and a union bound give total failure probability
at most $pB e^{-(k+1)/4}\le pBe^{-k/4}$.

For the frozen statistic,
\[
 \check H_{jb}=\xi_j(u_b)L_{jb}+R_{jb},
 \qquad
 |R_{jb}|\le b_{j,b}L_{jb}.
\]
The exponential Chernoff bound gives
$\Pp(L_{jb}>2\mid\bm U_j)\le2e^{-k/4}$.  Hence, outside an additional event
of probability at most $2pBe^{-k/4}$,
\[
 \frac1B\sum_{b=1}^B|R_{jb}|
 \le\frac2B\sum_{b=1}^Bb_{j,b}\le2\bar b_n
\]
for every $j$.

It remains to control
\[
 Z_j=\frac1B\sum_{b=1}^B
 \xi_j(u_b)(L_{jb}-1).
\]
Apply Lemma~\ref{lem:weighted-exp} conditionally with
$a_b=\xi_j(u_b)/B\le\Gamma/B$, then take a union bound over $j$.  With
conditional, and hence unconditional, probability at least $1-2pe^{-x}$,
\[
 \max_{j\le p}|Z_j|
 \le\Gamma\left(\sqrt{\frac{2x}{Bk}}+\frac{x}{Bk}\right).
\]
Combining the spatial, tail, stochastic, and quadrature terms proves
\eqref{eq:block-finite-bound}.
\end{proof}

\begin{corollary}[Block-score rate and screening]
\label{cor:block-rate}
Suppose
\[
 m_n\asymp nh_n,
 \qquad B_n\asymp h_n^{-1},
 \qquad k_n\asymp n\alpha_nh_n,
\]
so that $B_nk_n\asymp n\alpha_n$.  Assume
\begin{align}
 \log(2p)&=o(nh_n^2),
 \label{eq:block-rank-cond}\\
 \log(p/h_n)&=o(n\alpha_nh_n),
 \label{eq:block-local-cond}\\
 \bar\omega_n(C h_n,\tau_n,\alpha_n)+\bar b_n+q_{n,B}
 &=o(\Delta_{\D}),
 \qquad \tau_n=(pn)^{-2},
 \label{eq:block-bias-cond}\\
 \log(2p)&=o(n\alpha_n\Delta_{\D}^2).
 \label{eq:block-screen-cond}
\end{align}
Then
\begin{equation}
 \max_{j\le p}|\widetilde\Psi_j-\Psi_j|
 =O_{\Pp}\!\left(
 \sqrt{\frac{\log(2p)}{n\alpha_n}}
 +\frac{\log(2p)}{n\alpha_n}
 \right)+o(\Delta_{\D}),
 \label{eq:block-rate-final}
\end{equation}
and every coordinate in $\D$ precedes every coordinate in $\D^c$ with
probability tending to one.  Hence the sure-screening, exact-recovery, and
baseline-threshold conclusions of the previous section hold with
$\widetilde\Psi$ in place of $\widehat\Psi$.
\end{corollary}

\begin{proof}
Choose $\eta_n$ exactly as in
Corollary~\ref{cor:asymptotic-identifiable};
\eqref{eq:block-rank-cond} makes the rank term vanish and ensures that the
spatial radius is at most $Ch_n$.  Since $B_n\asymp h_n^{-1}$,
\eqref{eq:block-local-cond} gives $pB_ne^{-k_n/4}\to0$.  The PIT truncation
term vanishes for $\tau_n=(pn)^{-2}$.  Finally choose
$x_n=\log(2p)+r_n$, where $r_n\to\infty$ sufficiently slowly.  The last term
in \eqref{eq:block-finite-bound} vanishes, and because
$B_nk_n\asymp n\alpha_n$, the stochastic radius has the order displayed in
\eqref{eq:block-rate-final}.  Conditions
\eqref{eq:block-bias-cond}--\eqref{eq:block-screen-cond} make the total error
$o_{\Pp}(\Delta_{\D})$, after which deterministic separation applies.
\end{proof}

\begin{proposition}[Block quadrature under bounded variation]
\label{prop:block-BV}
If $0\le\xi_j\le\Gamma$ and
$\max_j\TV_{\I}(\xi_j)\le V_n$, then
\begin{equation}
 q_{n,B}\le C_\varepsilon(V_n+\Gamma)
 \left(\frac{m_n}{n}+\frac1n\right).
 \label{eq:block-BV-bound}
\end{equation}
In the usual choice $m_n\asymp nh_n$, the quadrature error is
$O\{(V_n+\Gamma)h_n\}$, even for discontinuous profiles.
\end{proposition}

\begin{proof}
Let $d_n=m_n/(n+1)$ and associate to every midpoint $u_b$ the interval of
length $d_n$ centred at $u_b$.  These intervals are disjoint and contiguous.
Their union differs from $\I$ only by the two grid-endpoint errors and by the
unused final group of fewer than $m_n$ ranks; its symmetric-difference length
is $O(d_n+\delta_n)$.  On each block interval $J_b$,
\[
 \int_{J_b}|\xi_j(u)-\xi_j(u_b)|\,\mathrm du
 \le d_n\TV_{J_b}(\xi_j).
\]
Summing gives at most $d_nV_n$.  Normalizing by an interval length bounded
away from zero, and controlling the omitted boundary set by $\Gamma$, yields
\eqref{eq:block-BV-bound}.
\end{proof}

\section{The primitive doubling condition can be replaced}

The manuscript's deficit-mass doubling assumption is a convenient route to a
$1/t$ bound, but it is stronger than the transfer proof needs.  Define, with
$D_{j,u}=1/\gamma-1/\xi_j(u)$,
\begin{align*}
 B_{j,u}(t)
 &=\int e^{-tD_{j,u}(v)}c\{\iota_j(u,v)\}K_j(u,\mathrm dv),\\
 M_{j,u}(t)
 &=\int D_{j,u}(v)e^{-tD_{j,u}(v)}
   c\{\iota_j(u,v)\}K_j(u,\mathrm dv),\\
 \chi_n(T)
 &=\max_{j\le p_n}\sup_{u\in I_{\varepsilon/2}}
   \sup_{t\ge T}\frac{M_{j,u}(t)}{B_{j,u}(t)}.
\end{align*}

\begin{description}
\item[(P3-L)] \emph{Uniform Laplace localization:}
$\lim_{T\to\infty}\sup_n\chi_n(T)=0$.
\end{description}

This condition says exactly that the exponentially tilted fibre law
concentrates near zero reciprocal-index deficit.  It imposes no doubling,
regular variation, density, or polynomial margin.

\begin{theorem}[Primitive transfer under uniform Laplace localization]
\label{thm:laplace-transfer}
Assume the geometry and Lipschitz clauses of (P1), but not atomlessness;
assume (P2), (P4), and (P3-L).  Let
$x_\alpha=\log\{(3\alpha_n)^{-1}\}$ and
$c_\gamma=\underline\gamma/2$.  Then, after a common onset,
\begin{equation}
 b_n(3\alpha_n)
 \le6\Gamma^2\chi_n(c_\gamma x_\alpha)
   +4\Gamma^2\overline R_1(c_\gamma x_\alpha).
 \label{eq:laplace-E2}
\end{equation}
The spatial transfer bound obtained from (P4) is unchanged, and the
conditional-quantile representation follows from
Lemma~\ref{lem:randomized-pit}; no atomlessness assumption is required.
Consequently, the score and screening theorems apply whenever
\begin{equation}
 \chi_n(c_\gamma x_\alpha)
 +\overline R_1(c_\gamma x_\alpha)
 =o(\Delta_{\D}),
 \label{eq:laplace-rate}
\end{equation}
together with the corresponding spatial and quadrature conditions.
The original doubling assumption implies (P3-L) with
$\chi_n(T)=O(T^{-1})$, but is not necessary.
\end{theorem}

\begin{proof}
Suppress $(j,u)$.  Write
\[
 \mathcal A(t)=\int e^{-tD(v)}c(v)\{1+r(t,v)\}K(\mathrm dv).
\]
Then the projected survival function factors exactly as
\[
 \overline F_j(e^t\mid u)=e^{-t/\xi_j(u)}\mathcal A(t).
\]
For large $t$, $R_0(t)\le1/2$, and differentiation gives
\[
 \mathcal A'(t)
 =-\int De^{-tD}c(1+r)\,\mathrm dK
   +\int e^{-tD}c\,\partial_tr\,\mathrm dK.
\]
Since
$(1-R_0)B\le\mathcal A\le(1+R_0)B$,
\begin{equation}
 \left|\frac{\mathcal A'(t)}{\mathcal A(t)}\right|
 \le3\frac{M(t)}{B(t)}+2R_1(t)
 \le3\chi_n(t)+2R_1(t)=:g_n(t).
 \label{eq:A-derivative-L}
\end{equation}
In particular $g_n(t)\to0$ uniformly.

It remains to justify a uniform lower linear bound for the log-quantile.
Because $B'(t)=-M(t)$,
\[
 -\{\log B(t)\}'=\frac{M(t)}{B(t)}.
\]
Given $\eta>0$, (P3-L) gives a common $T_\eta$ such that this derivative is
at most $\eta$ for $t\ge T_\eta$.  Also
$B(T_\eta)\ge c_-e^{-T_\eta/\underline\gamma}$ uniformly.  Therefore
\[
 -\log B(t)\le C_\eta+\eta t,
 \qquad t\ge T_\eta,
\]
and hence $|\log\mathcal A(t)|=o(t)$ uniformly.  If
$T_{j,u}(x)=\log Q_j(e^x,u)$, quantile inversion gives
\[
 x=\frac{T_{j,u}(x)}{\xi_j(u)}
   -\log\mathcal A\{T_{j,u}(x)\}.
\]
For all sufficiently large $x$, the $o(T)$ term is bounded in absolute value
by $T/\underline\gamma$, so
$T_{j,u}(x)\ge c_\gamma x$ uniformly.

Let $a(t)=\mathcal A'(t)/\mathcal A(t)$.  Once
$\Gamma|a(t)|\le1/2$, differentiation of the inverse relation yields
\[
 \frac{\mathrm d}{\mathrm dx}\log\ell_j(e^x,u)
 =\frac{\xi_j(u)^2a\{T_{j,u}(x)\}}
        {1-\xi_j(u)a\{T_{j,u}(x)\}}.
\]
Using \eqref{eq:A-derivative-L} and
$T_{j,u}(x)\ge c_\gamma x$ gives
\[
 \sup_{j,u}\left|
 \frac{\mathrm d}{\mathrm dx}\log\ell_j(e^x,u)
 \right|
 \le6\Gamma^2\chi_n(c_\gamma x)
   +4\Gamma^2\overline R_1(c_\gamma x).
\]
The normalized increment defining $b_n(3\alpha_n)$ is an average of this
derivative over an interval starting at $x\ge x_\alpha$, proving
\eqref{eq:laplace-E2}.  The spatial argument uses only (P4) and the inverse
slope bound, so it is unchanged.  Finally, the original doubling condition
gives $M(t)/B(t)=O(1/t)$ by the tilted-deficit estimate, hence implies
(P3-L).
\end{proof}

\begin{proposition}[A still weaker near-max mass criterion]
\label{prop:nearmax-mass}
Let
\[
 H_{j,u}(z)=\int\mathbf1_{\{D_{j,u}(v)\le z\}}
 c\{\iota_j(u,v)\}K_j(u,\mathrm dv),
 \qquad
 h_n(z)=\inf_{j,u}H_{j,u}(z).
\]
For every $z>0$ and $T>0$,
\begin{equation}
 \chi_n(T)
 \le z+\frac{D_{\max}c_+e^{-Tz/2}}{h_n(z/2)},
 \qquad D_{\max}=\underline\gamma^{-1}.
 \label{eq:nearmax-chi}
\end{equation}
In particular, if for every fixed $z>0$,
$\inf_n h_n(z)>0$, then (P3-L) holds.  This requires only a uniformly
positive amount of fibre mass in every fixed neighbourhood of the maximal
tail index; it imposes no local shape as $z\downarrow0$.
\end{proposition}

\begin{proof}
Again suppress $(j,u)$.  Split the tilted first moment at $z$:
\[
 M(T)
 \le zB(T)+D_{\max}c_+e^{-Tz}.
\]
On the other hand,
\[
 B(T)\ge e^{-Tz/2}H(z/2).
\]
Dividing proves \eqref{eq:nearmax-chi}.  First choose $z$ small and then let
$T\to\infty$ to obtain uniform Laplace localization under the stated fixed
neighbourhood-mass condition.
\end{proof}


\section{An information-theoretic $1/\log n$ lower bound under (C1)--(C2)}

The logarithmic scale in the manuscript is not merely an artefact of the Hill
proof when the primitive class has no uniform tail onset.  The next result
uses the full data, so it is stronger than a lower bound for marginal
procedures.

\begin{theorem}[Delayed-onset non-identifiability]
\label{thm:delayed-onset}
Fix $\gamma_0>0$ and $a>0$.  There are two triangular arrays of models
$P_{0,n}$ and $P_{1,n}$ with $p=2$ which both satisfy (C1)--(C2), full fibre
support, atomlessness, and uniform bounds
$0<\underline\gamma\le\gamma\le\Gamma$ and
$0<c_-\le c(u)\le c_+<\infty$, such that:
\begin{enumerate}
\item coordinate $1$ is inactive under $P_{0,n}$ and detectable under
      $P_{1,n}$;
\item under $P_{1,n}$ its population gap satisfies
      \begin{equation}
       \Delta_{1,n}=\frac{a}{4\gamma_0\log n};
       \label{eq:lower-gap}
      \end{equation}
\item the total variation distance between the full $n$-sample experiments
      satisfies
      \begin{equation}
       \|P_{0,n}^{\otimes n}-P_{1,n}^{\otimes n}\|_{\mathrm{TV}}
       \le n^{-1}.
       \label{eq:lower-TV}
      \end{equation}
\end{enumerate}
Consequently, for every screening rule $\widehat\A_n$ based on the complete
sample $(U_i,Y_i)_{i=1}^n$,
\begin{equation}
 P_{0,n}(1\in\widehat\A_n)
 +P_{1,n}(1\notin\widehat\A_n)
 \ge1-n^{-1}.
 \label{eq:lower-risk}
\end{equation}
Thus no procedure can uniformly distinguish an inactive coordinate from a
detectable coordinate with gap of order $1/\log n$ over the unquantified
(C1)--(C2) class.
\end{theorem}

\begin{proof}
Let $U_1,U_2$ be independent standard uniforms under both models and set
\[
 T_n=n^{2\gamma_0},
 \qquad
 \delta_n=\frac{a}{\log T_n}=\frac{a}{2\gamma_0\log n}.
\]
Under $P_{0,n}$, let $\gamma_{0,n}(u_1,u_2)=\gamma_0$ and let
\[
 \overline F_{0,n}(y\mid u)=y^{-1/\gamma_0},
 \qquad y\ge1.
\]
Under $P_{1,n}$, let
\[
 \gamma_{1,n}(u_1,u_2)=\gamma_0+\delta_nu_1
\]
and define the conditional survival function by the delayed splice
\begin{equation}
 \overline F_{1,n}(y\mid u)=
 \begin{cases}
  y^{-1/\gamma_0}, & 1\le y\le T_n,\\[2mm]
  T_n^{-1/\gamma_0}(y/T_n)^{-1/\gamma_{1,n}(u)},
      & y>T_n.
 \end{cases}
 \label{eq:delayed-splice}
\end{equation}
This is continuous and decreasing.  For $y>T_n$ it can be written as
\[
 \overline F_{1,n}(y\mid u)
 =c_n(u)y^{-1/\gamma_{1,n}(u)},
 \qquad
 c_n(u)=T_n^{1/\gamma_{1,n}(u)-1/\gamma_0}.
\]
Moreover $c_n(u)\le1$ and
\[
 \log c_n(u)
 =-\frac{\delta_nu_1\log T_n}
        {\gamma_0\{\gamma_0+\delta_nu_1\}}
 \ge-\frac{a}{\gamma_0^2},
\]
so $e^{-a/\gamma_0^2}\le c_n(u)\le1$ uniformly.  The slowly varying factor
is exactly equal to $c_n(u)$ beyond $T_n$, hence (C1)--(C2) hold.  For large
$n$, the tail indices also obey any fixed positive lower and upper bounds
containing $\gamma_0$.

Under $P_{0,n}$ coordinate $1$ is inactive.  Under $P_{1,n}$ it is active and,
by full support,
\[
 \xi_{1,n}(u)=\gamma_0+\delta_nu,
 \qquad
 \gmax_n=\gamma_0+\delta_n.
\]
Since the average of $u$ on $\I=[\varepsilon,1-\varepsilon]$ is $1/2$,
\[
 \Delta_{1,n}
 =\delta_n\left(1-\frac12\right)
 =\frac{a}{4\gamma_0\log n},
\]
which is \eqref{eq:lower-gap}.

The two one-observation laws agree exactly on the event $\{Y\le T_n\}$,
and under both models
\[
 \Pp(Y>T_n\mid U)=T_n^{-1/\gamma_0}=n^{-2}.
\]
Therefore the two product measures agree on
$E_n=\{\max_{1\le i\le n}Y_i\le T_n\}$, which has the same probability under
both, and
\[
 \|P_{0,n}^{\otimes n}-P_{1,n}^{\otimes n}\|_{\mathrm{TV}}
 \le P_{0,n}(E_n^c)
 \le n\,n^{-2}=n^{-1}.
\]
This proves \eqref{eq:lower-TV}.  Applying the elementary testing inequality
$P_0(A)+P_1(A^c)\ge1-\|P_0-P_1\|_{\mathrm{TV}}$ to
$A=\{1\in\widehat\A_n\}$ gives \eqref{eq:lower-risk}.
\end{proof}

\begin{remark}[What the lower bound does and does not settle]
Theorem~\ref{thm:delayed-onset} shows that the order $1/\log n$ is
information-theoretically inaccessible uniformly over (C1)--(C2), because
those conditions permit the asymptotic tail regime to start after the sample
has effectively ended.  It does not rule out detecting such gaps under a
common quantitative onset condition such as (P2), or after an explicit
second-order bias correction.  It identifies precisely why a quantified
primitive tail condition is indispensable for any uniform theorem.
\end{remark}

\section{A sharp impossibility result for all coordinatewise marginal screens}

Flatness of the projected tail index is not merely a weakness of the average
score.  In general, even the entire marginal law of $(U_j,Y)$ can be
identical under a model where $j$ is inactive and a model where it is active.

\begin{theorem}[Exact marginal non-identifiability]
\label{thm:marginal-impossibility}
There exist two models satisfying (C1)--(C2), full fibre support, continuity of
$\gamma$, and exact conditional Pareto tails, such that coordinate $1$ is
inactive in the first model and active in the second, while the law of
$(U_1,Y)$ is exactly the same in both models.  Consequently, for every
possibly randomized decision rule
$\phi_n=\phi_n\{(U_{i1},Y_i)_{i=1}^n\}\in\{0,1\}$,
\begin{equation}
 \Pp_0(\phi_n=1)+\Pp_1(\phi_n=0)=1,
 \label{eq:impossibility-risk}
\end{equation}
so its worst-case error over the two models is at least $1/2$ for every $n$.
\end{theorem}

\begin{proof}
Let $U_1,U_2$ be independent standard uniforms and choose
\[
 g(t)=\gamma_0+a\cos(2\pi t),
 \qquad 0<a<\gamma_0.
\]
Under model $P_0$, set $\gamma_0(u_1,u_2)=g(u_2)$.  Under model $P_1$, set
\[
 \gamma_1(u_1,u_2)=g\{(u_1+u_2)\bmod1\}.
\]
Because $g$ is periodic, $\gamma_1$ is continuous on the square.  Generate
$Y$ conditionally by
\[
 \Pp_r(Y>y\mid U_1=u_1,U_2=u_2)
 =y^{-1/\gamma_r(u_1,u_2)},
 \qquad y\ge1,
 \quad r\in\{0,1\}.
\]
All stated regularity conditions hold with $c\equiv1$.  Coordinate $1$ is
inactive under $P_0$ and active under $P_1$.

For every $u_1$ and $y\ge1$,
\begin{align*}
 \Pp_1(Y>y\mid U_1=u_1)
 &=\int_0^1y^{-1/g\{(u_1+v)\bmod1\}}\,\mathrm dv\\
 &=\int_0^1y^{-1/g(v)}\,\mathrm dv
 =\Pp_0(Y>y\mid U_1=u_1),
\end{align*}
where the middle identity is invariance of Lebesgue measure under rotation of
the circle.  Hence the joint laws of $(U_1,Y)$ coincide.  The distribution of
$\phi_n$ is therefore the same under $P_0$ and $P_1$; if its common
probability of outputting one is $q$, the two errors are $q$ and $1-q$, whose
sum is one.
\end{proof}

\section{Group order: the exact price of escaping marginal invisibility}

For $J\subseteq\{1,\ldots,p\}$ define, under full support,
\[
 \xi_J(u_J)=\max_{v\in[0,1]^{p-|J|}}\gamma(u_J,v),
 \qquad
 \Delta_J=\gmax-\int_{[0,1]^{|J|}}\xi_J(u_J)\,\mathrm du_J.
\]

\begin{theorem}[Group detectability and minimal interaction order]
\label{thm:group-order}
Assume full fibre support and continuity of $\gamma$.  For every nonempty
$J$,
\begin{equation}
 \Delta_J>0
 \quad\Longleftrightarrow\quad
 \lambda_{|J|}\bigl([0,1]^{|J|}\setminus\pi_J(\mathcal M)\bigr)>0.
 \label{eq:group-geometry}
\end{equation}
For an active coordinate $j$, define
\[
 r_\star(j)=\min\bigl\{|J|:j\in J,\ \Delta_J>0\bigr\}.
\]
If $\A$ is the active set and $s=|\A|$, then $r_\star(j)\le s$ for every
$j\in\A$.  The bound is sharp: for
\[
 \gamma(u)=g\{(u_1+\cdots+u_s)\bmod1\},
\]
with nonconstant positive periodic continuous $g$, one has
$r_\star(j)=s$ for every $j\le s$.
\end{theorem}

\begin{proof}
The proof of \eqref{eq:group-geometry} is identical to the one-coordinate
argument: $\xi_J(u_J)=\gmax$ exactly when the $J$-fibre through $u_J$
intersects $\mathcal M$, equivalently when $u_J\in\pi_J(\mathcal M)$; the
integrated nonnegative deficit is positive exactly when the complement has
positive measure.

Take $J=\A$.  Then $\xi_\A(u_\A)=\gamma(u_\A)$ because no active coordinate
is maximized out.  Since $j\in\A$, the continuous function $\gamma$ is not
constant on $[0,1]^\A$, so the open set on which
$\gamma<\gmax$ is nonempty and has positive measure.  Thus
$\Delta_\A>0$ and $r_\star(j)\le s$.

For the cyclic-sum example, fix any $J$ with $|J|<s$ and any value $u_J$.
At least one active coordinate is omitted.  Choose that omitted coordinate,
modulo one, so that the total sum is a maximizer of $g$; set the other omitted
coordinates arbitrarily.  Hence every $u_J$ belongs to
$\pi_J(\mathcal M)$, so $\Delta_J=0$.  For $J=\A$, the preceding paragraph
gives $\Delta_\A>0$.  Therefore $r_\star(j)=s$.
\end{proof}

\section{Summary of the strengthened theorem chain}

The exact population target is $\D$, the positive-gap set generated by the
conditional supports.  Under null argmax accessibility, $\D\subseteq\A$; it
equals $\A$ precisely when every active coordinate misses the accessible
argmax on a positive-measure set.  Continuity of projected conditional laws
is unnecessary.  The original empirical-rank score consistently recovers
$\D$ under a direct quadrature condition, with the rank-displacement
requirement weakened from $\log(pn)=o(nh^2)$ to
$\log p=o(nh^2)$.  The disjoint-block score further improves the stochastic
score rate from
\[
 \sqrt{\frac{\log(pn)}{n\alpha h}}
 \quad\text{to}\quad
 \sqrt{\frac{\log p}{n\alpha}},
\]
while replacing uniform-in-location tail and spatial assumptions by their
block averages.  Finally, primitive deficit-mass doubling can be replaced by
the strictly weaker uniform Laplace-localization condition, and the
impossibility theorem shows that no purely coordinatewise marginal method can
remove the remaining argmax-projection limitation.

\end{document}